% Источники: exam/materials/моделирование.pdf, раздел 5; exam/materials/преподаватель/2020/01-05-2020-Model_L11__1_.pdf; exam/materials/преподаватель/2020/01-05-2020-Model_L11__2_.pdf; exam/materials/прошлогодние ответы/Modelirovanie_1_-3.pdf, вопрос 35; GitHub HumsterProgrammer/iu9-notes/8sem/Моделирование/Моделирование_лекция-15_02-04-26.md.
\section{Планирование эксперимента. Факторные планы, дробные факторные планы.}

\textbf{Объект в теории планирования эксперимента} рассматривается как «черный ящик», имеющий входы $v$ (управляемые независимые параметры) и выходы $y$. Переменные $v$ принято называть факторами.

\textbf{Активный эксперимент} включает: систему воздействий, при которых воспроизводится функционирование объекта; регистрацию отклика объекта. Для него строится план эксперимента.

Предполагается активный тип экспериментов, когда имеется возможность независимо и целенаправленно менять значения факторов $v$ во всем требуемом диапазоне. Факторы бывают качественными и количественными. Качественные факторы можно квантифицировать или приписать им числовые обозначения, тем самым перейти к количественным значениям.

Переменным $v$ можно сопоставить геометрическое понятие факторного пространства — пространства, координатные оси которого соответствуют значениям факторов. Совокупность конкретных значений всех факторов образует точку в многомерном факторном пространстве. На объект или систему могут влиять возмущающие факторы, но они являются случайными и не поддаются управлению.

Область планирования задается интервалами возможного изменения факторов:
\[
  v_{i,\min} < v_i < v_{i,\max}, \qquad i=1,2,\ldots,k,
\]
где $k$ — количество факторов.

В теории планирования экспериментов часто используют нормализацию факторов, то есть преобразование натуральных значений факторов в безразмерные (кодированные) величины:
\[
  x_i = \frac{v_i-v_{i0}}{\Delta v_i},
\]
где $v_i$ — натуральное значение фактора, $v_{i0}$ — натуральное значение основного уровня фактора, соответствующее нулю в безразмерной шкале, $\Delta v_i$ — интервал варьирования.

Совокупность основных уровней всех факторов представляет собой точку в пространстве параметров, называемую центральной точкой плана или центром эксперимента.

\textbf{План эксперимента} — совокупность данных, определяющих число, условия и порядок реализации опытов.

\textbf{Опыт} — воспроизведение исследуемого явления в конкретных условиях с последующей регистрацией результата. В условиях случайности в одних и тех же условиях проводятся параллельные (повторные) опыты в интересах получения статистически устойчивых результатов.

\textbf{Матрица плана} — стандартная форма записи условий проведения экспериментов в виде прямоугольной таблицы, строки которой отвечают опытам, столбцы — факторам (последний столбец — отклик — наблюдаемая случайная переменная, по предположению, зависящая от факторов — его обычно нет).

\textbf{Полный факторный план} — план, содержащий все возможные комбинации всех факторов на определенном числе уровней равное число раз.

\textbf{Дробный факторный план} — план, содержащий часть комбинаций полного факторного плана.

На первом этапе планирования эксперимента необходимо выбрать область определения факторов $x_i$. Выбор этой области производится исходя из априорной информации. Значения $x_i$ называются уровнями управляющего параметра.

\clearpage
